\documentclass[10pt,a4paper]{amsart}
\usepackage[margin=19mm]{geometry}
\usepackage{amsmath,amssymb,mathtools}
\usepackage[scr=rsfs,cal=euler]{mathalfa}
\usepackage{xfrac}
\usepackage[unicode,hidelinks,bookmarks=false]{hyperref}
\newcommand{\ie}{i.e.\ }
\newcommand{\cf}{cf.\ }
\newcommand{\resp}{respectively\ }
\newcommand{\Subsection}{\subsection}
\providecommand{\nobreakdash}{\nobreak\hbox{-}}
\setlength{\emergencystretch}{2em}
\multlinegap0pt
\usepackage{float}
\usepackage{tikz}
\usepackage{amssymb}
\newtheorem{theorem}{Theorem}
\DeclareMathOperator{\Imag}{Im}
\DeclareMathOperator{\Real}{Re}
\newcommand{\cX}{\mathcal{X}}
\newcommand{\iu}{\mathrm{i}} % imaginary unit
\title{A perfectly matched layer approach \\ for the spectral split-step Pad\'e method}
\author{Daniel Walsken, Matthias Ehrhardt, Pavel Petrov}
\date{Source: arXiv:2607.03521v1 (CC BY 4.0)}
\begin{document}
\maketitle

\noindent\emph{Source and licence.} A perfectly matched layer approach \\ for the spectral split-step Pad\'e method. Version arXiv:2607.03521v1 (CC BY 4.0), distributed by arXiv under the Creative Commons Attribution 4.0 International licence, http://creativecommons.org/licenses/by/4.0/, as recorded in the arXiv licence metadata for this exact version and shown on the abstract page https://arxiv.org/abs/2607.03521v1. Full licence notice: \texttt{licenses/arxiv-2607.03521-CC-BY-4.0.txt}.\par\smallskip

\noindent\textit{Editorial scope.} Counted unit: the article's theorem theorem (source0.tex lines 808--825). The article's complete original proof of it is delivered with it (source0.tex lines 827--864, one \texttt{proof} environment of 38 lines). No mathematics is written, repaired or re-derived: the statement and the proof are the article's own lines, in the article's own notation, and no retained line is substituted or re-flowed.  No further source-local context is needed: every cross-reference of every retained line resolves inside the two delivered spans.  Nothing was omitted from any retained span, so the package carries no omission record.  No retained line contains a citation, so the wrapper carries no bibliography and no bibliography entry.  No retained line contains a figure environment, an image-inclusion command or drawing code, so no image is delivered and the package declares no asset dependency.  Editorial formatting only, in addition: an AMS \texttt{amsart} preamble that restates the article's own declarations and macros used by the retained lines, namely the environments \texttt{theorem} on separate counters, together with the standard packages those lines need.  The retained environments print in this document as Theorem 1 in order of appearance, whereas the article prints the same items with the numbering of its own full document.  The complete native source of the article as distributed in the arXiv e-print for this exact version is bundled unchanged as \texttt{sources/204398/source0.tex}.
\begin{theorem}[Discrete Fourier stability of the SSSP--PML scheme]
We assume:
\begin{itemize}
\item $\Real S_y(y) \ge 1$, $\Imag S_y(y) \ge 0$,
\item $\Real b_j > 0$ for all Pad\'e coefficients,
\item the function $k^2(y)$ is bounded.
\end{itemize}
Then for each Fourier mode $\kappa_p$, the corresponding amplification factor $G_p$ of one SSSP step satisfies
\begin{equation*}
     |G_p| \le 1 + C h,
\end{equation*}
with a constant $C$ independent of $p$ and $N$. 
%In particular, 
For sufficiently small $h$, the scheme is stable in the $\ell^2$-sense
\begin{equation*}
\|u^{n+1}_N\|_{\ell^2} \le (1 + C h)\|u^n_N\|_{\ell^2}.
\end{equation*}
\end{theorem}

\begin{proof}
For each Fourier mode, the operator $\cX_{\mathrm{PML}}$ acts approximately as multiplication by
\begin{equation*}
     \lambda_p(y) = \frac{k^2(y) - k_0^2 - \kappa_p^2 / S_y^2(y)}{k_0^2}.
\end{equation*}
Since $\Real S_y \ge 1$ and $\Imag S_y \ge 0$, we obtain
\begin{equation*}
   \Real\Bigl(-\frac{\kappa_p^2}{S_y^2}\Bigr) \le 0,
\end{equation*}
which implies
%\begin{equation*}
    $\Real \lambda_p \le C$.
%\end{equation*}
%
Next, the Pad\'e approximation yields the amplification factor
\begin{equation*} 
   G_p = \mathrm{e}^{\iu k_0 h} \biggl(d_0 + \sum_{j=1}^P \frac{d_j}{1 + b_j \lambda_p} \biggr).
\end{equation*}
Since $\Real b_j > 0$ and $\Real \lambda_p \le C$, the denominator satisfies
%\begin{equation*}
$|1 + b_j \lambda_p| \ge c > 0$,
%\end{equation*}
uniformly in $p$. Hence
\begin{equation*}
\Bigl| \frac{d_j}{1 + b_j \lambda_p} \Bigr| \le C_j.
\end{equation*}
Now, summing over $j$ gives
%\begin{equation*}
$|G_p| \le 1 + C h$,
%\end{equation*}
where we used the consistency of the Pad\'e approximation.
%
Finally, Parseval's identity yields
\begin{equation*}
   \|u^{n+1}_N\|_{\ell^2}^2 = \sum_p |G_p|^2 |\hat{u}_p|^2 \le (1 + C h)^2 \|u^n_N\|_{\ell^2}^2,
\end{equation*}
which proves the result.
\end{proof}

\end{document}
